\section*{Chapter Summary}
This chapter detailed the technical implementation of the DBaaS platform, bridging the gap between high-level architectural design (Chapter~\ref{chap:system_design}) and the underlying Kubernetes runtime. By utilizing Go and the Kubernetes client ecosystem, the backend separates the control-plane REST API from data-plane components such as tenant databases, PostgREST resources, and PGProxy. Delegating database lifecycle management to CloudNativePG and the MongoDB Community Operator keeps provisioning declarative, while the backend focuses on validation, metadata, status reconciliation, connection management, and query execution.

The platform's presentation layer, built with Flutter, successfully abstracts this infrastructure complexity. As demonstrated in this chapter, the implemented user interfaces—ranging from SQL and NoSQL query workbenches to comprehensive management dashboards—allow users to provision, manage, and monitor database health without direct interaction with Kubernetes manifests. Security is rigorously maintained through Argon2id authentication, JWT token rotation, and role-based access control, coupled with a validation layer that prevents high-risk database operations.

Finally, the performance evaluation on the local k3d cluster measured provisioning latency and resource overhead, providing a baseline for future scalability. While the current deployment relies on containerized Go binaries and local k3d automation, the architecture is designed for production-grade hardening, including managed metadata services, stricter network policies, and multi-replica deployments.